%! Author = nursace
%! Date = 15.09.2026

\documentclass[aspectratio=169,11pt]{beamer}
\usepackage{color}
% Lecture 01 — Introduction to Java
% Course: Introduction to Programming with Java

\input{preamble}

\title{Lecture 01: Introduction to Java}
\subtitle{Introduction to Programming with Java}

\begin{document}

%----------------------------------------------------------------------
    \begin{frame}
        \titlepage
    \end{frame}

%----------------------------------------------------------------------
    \begin{frame}{Agenda}
        \begin{enumerate}
            \item What is Java?
            \item A short history
            \item Why learn Java?
            \item How Java runs: source $\rightarrow$ bytecode $\rightarrow$ JVM
            \item Where Java is used
            \item How this course approaches Java
            \item Looking ahead
        \end{enumerate}
    \end{frame}

%======================================================================
    \section{What is Java?}
%======================================================================

    \begin{frame}{What is Java?}
        \textbf{Java} is a general-purpose, object-oriented programming language.
        \vspace{0.8em}
        \begin{itemize}
            \item Designed for clarity, safety, and portability
            \item Programs are written once and run on many platforms
            (via the Java Virtual Machine)
            \item Large standard library and huge ecosystem of tools and frameworks
            \item Used in education, industry, research, and open-source software
        \end{itemize}
    \end{frame}

    \begin{frame}{Java is \ldots}
        \begin{columns}[T]
            \begin{column}{0.48\textwidth}
                \textbf{A language}
                \begin{itemize}
                    \item Syntax and rules
                    \item Types, classes, methods
                    \item What you type in \javacode{.java} files
                \end{itemize}
                \vspace{1em}
                \textbf{A platform}
                \begin{itemize}
                    \item JVM + libraries
                    \item Tools (compiler, debugger)
                    \item Runtime services
                \end{itemize}
            \end{column}
            \begin{column}{0.48\textwidth}
                \textbf{An ecosystem}
                \begin{itemize}
                    \item Build tools (Maven, Gradle)
                    \item IDEs (IntelliJ, Eclipse, VS~Code)
                    \item Frameworks (Spring, Jakarta EE, \ldots)
                    \item Community and documentation
                \end{itemize}
            \end{column}
        \end{columns}
    \end{frame}

    \begin{frame}{A first mental model}
        \begin{block}{Three layers (for now)}
            \begin{enumerate}
                \item \textbf{You} write Java source code (\javacode{.java})
                \item The \textbf{compiler} (\javacode{javac}) turns it into bytecode (\javacode{.class})
                \item The \textbf{JVM} executes that bytecode on your machine
            \end{enumerate}
        \end{block}
        \vspace{0.6em}
        You do not need hardware-specific machine code for each OS---
        the JVM handles that difference.
    \end{frame}

    \begin{frame}[fragile]{Hello, Java (preview)}
        \begin{lstlisting}
public class Hello {
    public static void main(String[] args) {
        System.out.println("Hello, Java!");
    }
}
        \end{lstlisting}
        \vspace{0.4em}
        \begin{itemize}
            \item Every application needs a \javacode{main} method as an entry point
            \item \javacode{System.out.println} prints a line to the console
            \item Class name and file name usually match: \javacode{Hello.java}
        \end{itemize}
    \end{frame}

    \begin{frame}{We will unpack this later}
        For now, treat \javacode{public class \ldots} and \javacode{public static void main}
        as the standard wrapper for a runnable program.
        \vspace{0.8em}
        \begin{itemize}
            \item Later lectures: classes, methods, \javacode{static}, access modifiers
            \item Today: motivation, platform model, and where Java fits
        \end{itemize}
        \vspace{0.6em}
        \trythis{If you already have a JDK installed, compile and run
        \javacode{Hello.java}. If not, we will set up the toolchain soon.}
    \end{frame}

    %--- Quiz checkpoint 1 (after What is Java?) ---
    \begin{frame}{Quiz 1 --- What is Java?}
        \quizbox{Which statement best describes Java?}
        \begin{enumerate}
            \item A markup language for web pages (like HTML)
            \item A general-purpose, object-oriented language that runs via a JVM
            \item A database query language
            \item An operating system for servers
        \end{enumerate}
        \quizpause
    \end{frame}

    \begin{frame}{Answer --- Quiz 1}
        \answerbox{\textbf{2.} A general-purpose, object-oriented language that runs via a JVM.}
        \vspace{0.4em}
        Java programs are written as source, compiled to bytecode, and executed by the
        Java Virtual Machine---not by the CPU directly as native machine code.
    \end{frame}

%======================================================================
    \section{A short history}
%======================================================================

    \begin{frame}{A short history of Java}
        \begin{itemize}
            \item \textbf{1991--1995:} Started at Sun Microsystems (project ``Oak'', then ``Java'')
            \item Goal: write software that could run on many devices and networks
            \item \textbf{1995:} Java 1.0 announced --- ``Write Once, Run Anywhere''
            \item Grew with the web (applets era), then enterprise servers
            \item \textbf{2010:} Oracle acquires Sun; Java continues as a major platform
            \item Today: long-term support (LTS) releases such as Java~17 and Java~21
        \end{itemize}
    \end{frame}

    \begin{frame}{Design ideas that still matter}
        From the early design, Java emphasised:
        \begin{itemize}
            \item \textbf{Simplicity} relative to C/C++ (no manual memory management for most code)
            \item \textbf{Object orientation} as a primary organising idea
            \item \textbf{Safety}: strong typing, array bounds checks, sandboxing ideas
            \item \textbf{Portability} via bytecode + JVM
            \item \textbf{Productivity}: rich standard library
        \end{itemize}
        \vspace{0.5em}
        These choices explain much of Java's classroom and industry success.
    \end{frame}

    \begin{frame}{Versions: what you need to know}
        \begin{itemize}
            \item Java has a versioned language \emph{and} platform
            \item New features arrive regularly; LTS versions are common in courses and companies
            \item This course uses \textbf{modern Java} (Java~17-friendly examples are fine)
            \item You rarely need to memorise version numbers early on
        \end{itemize}
        \vspace{0.6em}
        \begin{block}{Practical tip}
            Check your version with \javacode{java -version}.
            Prefer a current LTS JDK for coursework.
        \end{block}
    \end{frame}

%======================================================================
    \section{Why learn Java?}
%======================================================================

    \begin{frame}{Career relevance}
        \begin{itemize}
            \item Java remains a top language in enterprise software
            \item Common in banking, telecom, government, e-commerce, and SaaS
            \item Android (historically and still widely) and JVM-based backends
            \item Strong hiring signal: types, OOP, testing, and large codebases
        \end{itemize}
        \vspace{0.6em}
        Learning Java is not only about one language---it trains habits used in
        professional development.
    \end{frame}

    \begin{frame}{Portability: Write Once, Run Anywhere}
        \begin{columns}[T]
            \begin{column}{0.55\textwidth}
                \begin{itemize}
                    \item Same \javacode{.class} files can run on Windows, macOS, Linux, servers, containers
                    \item Provided a compatible JVM is installed
                    \item Reduces ``it works on my machine'' friction \emph{at the language/runtime level}
                \end{itemize}
            \end{column}
            \begin{column}{0.40\textwidth}
                \begin{block}{Caveat}
                    Portability is not magic: file paths, encodings, and native libraries still need care.
                \end{block}
            \end{column}
        \end{columns}
    \end{frame}

    \begin{frame}{Ecosystem and tooling}
        Java's ecosystem helps you grow beyond toy programs:
        \begin{itemize}
            \item Standard library: collections, I/O, concurrency, networking, \ldots
            \item Build and dependency tools (Maven, Gradle)
            \item Mature IDEs with refactoring and debugging
            \item Testing frameworks (e.g.\ JUnit) used in industry
            \item Huge amount of tutorials, books, and Stack Overflow answers
        \end{itemize}
    \end{frame}

    \begin{frame}{OOP foundation}
        Java is a natural vehicle for learning \textbf{object-oriented programming}:
        \begin{itemize}
            \item Classes and objects as everyday building blocks
            \item Encapsulation, inheritance, and interfaces (later in the course)
            \item Clear separation of ``what an object knows'' and ``what it can do''
        \end{itemize}
        \vspace{0.6em}
        These ideas transfer to C\#, Kotlin, TypeScript, and many other languages.
    \end{frame}

    \begin{frame}{Why Java in a first-year course?}
        \begin{itemize}
            \item Forces you to be explicit about types and structure
            \item Errors often appear early (compile-time), which helps learning
            \item Scales from small exercises to multi-file projects
            \item Matches many follow-on courses (algorithms, software engineering, databases)
        \end{itemize}
        \vspace{0.5em}
        \commonmistake{Thinking Java is ``only for big companies.''
        It is also excellent for learning core CS ideas carefully.}
    \end{frame}

    %--- Quiz checkpoint 2 (after Why learn Java?) ---
    \begin{frame}{Quiz 2 --- Why learn Java?}
        \quizbox{True or false?\\[0.4em]
        ``Write Once, Run Anywhere'' means a Java \javacode{.class} file can run on
        different operating systems \emph{without} rewriting the program, as long as
        a compatible JVM is available.}
        \begin{enumerate}
            \item True
            \item False
        \end{enumerate}
        \quizpause
    \end{frame}

    \begin{frame}{Answer --- Quiz 2}
        \answerbox{\textbf{True.} Bytecode is portable across platforms that provide a compatible JVM.}
        \vspace{0.4em}
        Caveat (from earlier): paths, encodings, and native libraries can still differ---
        portability is strong at the language/runtime level, not magic for every detail.
    \end{frame}

%======================================================================
    \section{How Java runs}
%======================================================================

    \begin{frame}{From idea to running program}
        High-level pipeline:
        \vspace{0.6em}
        \begin{center}
            \begin{tikzpicture}[
                box/.style={draw=LectureNavy, rounded corners, align=center,
                fill=CodeBg, minimum height=1.1cm, minimum width=2.4cm, font=\small},
                arr/.style={-{Latex}, thick, LectureAccent}
            ]
                \node[box] (src) {Source\\{\ttfamily .java}};
                \node[box, right=0.7cm of src] (cmp) {Compiler\\{\ttfamily javac}};
                \node[box, right=0.7cm of cmp] (bc) {Bytecode\\{\ttfamily .class}};
                \node[box, right=0.7cm of bc] (jvm) {JVM};
                \node[box, right=0.7cm of jvm] (out) {Output};
                \draw[arr] (src) -- (cmp);
                \draw[arr] (cmp) -- (bc);
                \draw[arr] (bc) -- (jvm);
                \draw[arr] (jvm) -- (out);
            \end{tikzpicture}
        \end{center}
        \vspace{0.5em}
        Keep this pipeline in mind whenever something ``does not run'':
        Was it compiled? Is the class name correct? Is a JDK available?
    \end{frame}

    \begin{frame}{Source code}
        \begin{itemize}
            \item Human-readable text in \javacode{.java} files
            \item Must follow Java's syntax rules
            \item Organised into classes (and later packages)
            \item Edited in any text editor or IDE
        \end{itemize}
        \vspace{0.5em}
        \begin{block}{Analogy}
            Source code is the recipe. The computer cannot ``cook'' it directly
            until it is compiled and executed.
        \end{block}
    \end{frame}

    \begin{frame}{Bytecode (high level)}
        \begin{itemize}
            \item The compiler produces \textbf{bytecode}, not native CPU instructions
            \item Bytecode is a compact instruction set for the JVM
            \item Stored typically in \javacode{.class} files
            \item Same bytecode can be shipped to different operating systems
        \end{itemize}
        \vspace{0.5em}
        You do not write bytecode by hand in this course.
    \end{frame}

    \begin{frame}{The Java Virtual Machine (JVM)}
        The \textbf{JVM} is a program that:
        \begin{itemize}
            \item Loads bytecode
            \item Verifies it (basic safety checks)
            \item Executes it (often with Just-In-Time compilation for speed)
            \item Manages memory (garbage collection)
        \end{itemize}
        \vspace{0.5em}
        Think of the JVM as a portable ``computer'' implemented in software.
    \end{frame}

    \begin{frame}{JDK vs JRE (student view)}
        \begin{itemize}
            \item \textbf{JDK} (Java Development Kit): compiler + tools + runtime --- what \emph{you} install to write programs
            \item \textbf{JRE} (Java Runtime Environment): enough to \emph{run} programs (concept historically important)
            \item In modern distributions, installing a JDK is the usual student choice
        \end{itemize}
        \vspace{0.5em}
        \begin{block}{Command-line essentials}
            \javacode{javac Hello.java} \quad then \quad \javacode{java Hello}
        \end{block}
    \end{frame}

    \begin{frame}[fragile]{Compile and run (workflow)}
        \begin{lstlisting}
# In a terminal (comments for humans):
javac Hello.java    # creates Hello.class
java Hello          # runs the main method
        \end{lstlisting}
        \vspace{0.4em}
        \begin{itemize}
            \item Do not include \javacode{.class} in the \javacode{java} command
            \item Class name is case-sensitive
        \end{itemize}
        \vspace{0.3em}
        \commonmistake{Running \javacode{java Hello.java} by habit from other tools.
        Classic workflow is \javacode{javac} then \javacode{java ClassName}.
            (Some newer JDKs can run single-file source directly---nice, but learn the classic steps.)}
    \end{frame}

    \begin{frame}{Why this model helps beginners}
        \begin{itemize}
            \item Compile-time errors catch many mistakes early
            \item Clear separation: editing vs compiling vs running
            \item Encourages reading error messages carefully
            \item Mirrors how larger projects are built
        \end{itemize}
    \end{frame}

    %--- Quiz checkpoint 3 (after How Java runs) ---
    \begin{frame}{Quiz 3 --- Compile and run}
        \quizbox{You have a file \javacode{Hello.java}. Which pair is the classic workflow?}
        \begin{enumerate}
            \item \javacode{java Hello.java} then \javacode{javac Hello}
            \item \javacode{javac Hello.java} then \javacode{java Hello}
            \item \javacode{javac Hello} then \javacode{java Hello.class}
            \item \javacode{java Hello} only (no compile step ever)
        \end{enumerate}
        \quizpause
    \end{frame}

    \begin{frame}{Answer --- Quiz 3}
        \answerbox{\textbf{2.} \javacode{javac Hello.java} creates \javacode{Hello.class}; then \javacode{java Hello} runs it.}
        \vspace{0.4em}
        Do not pass \javacode{.class} to the \javacode{java} command. Some newer JDKs can run
        a single source file directly---useful later, but learn the classic two steps first.
    \end{frame}

    \begin{frame}{Quiz 3b --- Bytecode (true/false)}
        \quizbox{True or false?\\[0.4em]
        The Java compiler produces native CPU machine code for your laptop's processor.}
        \begin{enumerate}
            \item True
            \item False
        \end{enumerate}
        \quizpause
    \end{frame}

    \begin{frame}{Answer --- Quiz 3b}
        \answerbox{\textbf{False.} The compiler produces \textbf{bytecode} for the JVM, not CPU-specific machine code.}
        \vspace{0.4em}
        The JVM (often with JIT compilation) handles execution on the actual hardware.
    \end{frame}

%======================================================================
    \section{Where Java is used}
%======================================================================

    \begin{frame}{Where Java is used}
        Java appears across many domains:
        \begin{itemize}
            \item \textbf{Backend / server} applications and APIs
            \item \textbf{Android} applications (Java/Kotlin on the JVM/Android runtime family)
            \item \textbf{Desktop} tools and cross-platform apps
            \item \textbf{Big data} and streaming (e.g.\ parts of Hadoop/Kafka ecosystems)
            \item \textbf{Scientific} and financial systems
            \item \textbf{Embedded} and IoT in selected environments
        \end{itemize}
    \end{frame}

    \begin{frame}{Enterprise and the web}
        \begin{itemize}
            \item Many companies run large Java services for years
            \item Frameworks help structure web apps and microservices
            \item Strong typing and tooling support large teams
            \item Reliability and maintainability often matter as much as raw speed
        \end{itemize}
        \vspace{0.5em}
        You will not build a full enterprise stack in week one---
        but the language you learn here is the same foundation.
    \end{frame}

    \begin{frame}{Related languages on the JVM}
        The JVM hosts more than Java:
        \begin{itemize}
            \item \textbf{Kotlin} --- popular for Android and modern JVM apps
            \item \textbf{Scala}, \textbf{Clojure}, and others for specialised styles
        \end{itemize}
        \vspace{0.5em}
        Learning Java first makes these easier to approach later.
        \vspace{0.5em}
        \trythis{After class, skim the Oracle or OpenJDK ``What is Java?'' overview page
        and note three facts you did not know.}
    \end{frame}

    \begin{frame}{Java vs other first languages (honest comparison)}
        \begin{tabular}{@{}lll@{}}
            \toprule
            \textbf{Aspect} & \textbf{Java} & \textbf{Often compared to} \\
            \midrule
            Typing & Static, explicit & Python (dynamic) \\
            Verbosity & More ceremony early & Scripting languages \\
            Feedback & Compile + run & Often interpret/run \\
            OOP & Central & Varies \\
            Industry & Very strong & Depends on domain \\
            \bottomrule
        \end{tabular}
        \vspace{0.6em}
        None is universally ``best''---Java is a strong teaching and career choice.
    \end{frame}

    %--- Quiz checkpoint 4 (after Where Java is used) ---
    \begin{frame}{Quiz 4 --- Where Java is used}
        \quizbox{Which domain is a \emph{typical} place you will meet Java in industry?}
        \begin{enumerate}
            \item Backend / server applications and APIs
            \item Replacing HTML as the only language for static web pages
            \item Writing GPU shader programs instead of GLSL
            \item Assembling BIOS firmware for every laptop brand
        \end{enumerate}
        \quizpause
    \end{frame}

    \begin{frame}{Answer --- Quiz 4}
        \answerbox{\textbf{1.} Backend / server applications and APIs (also Android/JVM apps, tools, finance, and more).}
        \vspace{0.4em}
        Java is not ``only enterprise,'' but large server-side systems are a very common home for it.
    \end{frame}

%======================================================================
    \section{How this course approaches Java}
%======================================================================

    \begin{frame}{Course philosophy}
        We will:
        \begin{itemize}
            \item Start with syntax and small programs you can reason about
            \item Emphasise reading code, not only writing it
            \item Introduce OOP concepts carefully, with examples
            \item Practise debugging and explaining errors
            \item Build toward structured, multi-class programs
        \end{itemize}
    \end{frame}

    \begin{frame}{Habits that help}
        \begin{itemize}
            \item Type examples yourself---do not only read slides
            \item Change one thing at a time and observe the result
            \item Keep programs small until ideas are solid
            \item Name variables meaningfully (\javacode{score}, not \javacode{x1})
            \item Ask: ``What is the state of my program after this line?''
        \end{itemize}
    \end{frame}

    \begin{frame}{Tools you will use}
        \begin{itemize}
            \item A JDK (OpenJDK or another distribution)
            \item A terminal for \javacode{javac} / \javacode{java}
            \item An IDE or editor (course recommendation is IntelliJ IDEA)
            \item Version control later (e.g.\ Git) for projects
        \end{itemize}
        \vspace{0.5em}
        Setup details belong in a lab session; conceptually you now know \emph{why}
        the JDK matters.
    \end{frame}

%======================================================================
    \section{Looking ahead}
%======================================================================

    \begin{frame}{Key takeaways}
        \begin{enumerate}
            \item Java is a portable, object-oriented language with a large ecosystem
            \item Career, portability, tools, and OOP foundations motivate learning it
            \item Source $\rightarrow$ bytecode $\rightarrow$ JVM is the core execution story
            \item Java powers servers, Android, tools, and many large systems
            \item This course builds from small, clear programs toward OOP design
        \end{enumerate}
    \end{frame}

    \begin{frame}{Questions?}
        \begin{center}
            {\Large What questions do you have about Java so far?}
        \end{center}
        \vspace{1.2em}
        \begin{block}{Suggested reading / practice}
            \begin{itemize}
                \item Skim an official ``Introduction to Java'' overview
                \item Install a JDK and run \javacode{java -version}
                \item Compile and run a \javacode{Hello} program
            \end{itemize}
        \end{block}
    \end{frame}

\end{document}
